%=====================================================================
%  CAD Design Package - Mechanical Drawing Set and Gazebo Model
%  Autonomous TEC-Based Atmospheric-Water PV-Cleaning Robot (REE598)
%
%  Companion to prototype_design_report.tex. All frozen numbers come
%  from scripts/tec_robot_design.py via output/tec_robot_results.json.
%  Drawings: cad/D-101..D-107 (.scr for AutoCAD + .dxf via ezdxf).
%  Simulation model: gazebo/pvbot_description/.
%  NOTE: colortbl deliberately not used (MiKTeX array-version bug).
%=====================================================================
\documentclass[11pt,a4paper]{article}

\usepackage[utf8]{inputenc}
\usepackage[T1]{fontenc}
\usepackage{lmodern}
\usepackage[english]{babel}
\usepackage{amsmath,amssymb}
\usepackage{graphicx}
\usepackage{booktabs}
\usepackage{array}
\usepackage{longtable}
\usepackage{geometry}
\geometry{margin=2.4cm}
\usepackage{cite}
\usepackage[colorlinks=true,linkcolor=blue,citecolor=blue,urlcolor=blue]{hyperref}

\newcommand{\dwg}[1]{\textbf{#1}}
\newcommand{\fpath}[1]{{\ttfamily\small #1}}

\title{\textbf{CAD Design Package}\\[0.4em]
\large Mechanical Drawing Set (D-101 to D-110), 3D Model, and Gazebo Model\\[0.2em]
\normalsize Autonomous Thermoelectric Atmospheric-Water PV-Cleaning Robot --- REE598}
\author{Rahma El-Hamouly}
\date{13 June 2026}

\begin{document}
\maketitle

\section{Purpose and scope}
This package turns the frozen design values of the prototype design report
into manufacturable 2D drawings, a 3D model, and a simulation model. The robot
has since been revised from a free differential-drive crawler to a rail-guided
carriage that carries a full panel-width brush (Section~\ref{sec:railrev}
gives the reason and the changed numbers). The drawing set covers every
structural part of the 8.0\,kg carriage: chassis plate, motor brackets, brush
assembly, the 4-tile TEC condenser sandwich, the 3.6\,L water tank with its
pump, the electronics bay, a general assembly, and the three rail-drive sheets
added by the revision (D-108 rail assembly, D-109 rack-and-pinion drive, D-110
full-width brush). Each 2D drawing exists in two formats generated from the
same dimensions: an AutoCAD script (\fpath{.scr}, plain text, no dependencies)
and a DXF file (\fpath{.dxf}, written with the \fpath{ezdxf} library by
\fpath{scripts/make\_cad\_dxf.py}). The 3D model is in \fpath{cad/3d/}:
AutoCAD scripts that build native 3D solids of the carriage and the panel/rail
scene, plus a faceted-mesh \fpath{.dxf} for any viewer. The Gazebo model in
\fpath{gazebo/pvbot\_description/} uses the hand-maintained URDF, with the same
link names and inertial values, so the mechanical and software documents
describe one robot.

The frozen inputs are: carriage mass 8.0\,kg; full-width brush 0.96\,m at
$\varnothing$50\,mm with microfibre roller and low-pressure spray; operation
on a residential panel tilted $24^\circ$; rail-drive force 37.2\,N at the
rack, needing at least 0.84\,N\,m at the $\varnothing$30\,mm pinion turning
near 38\,rpm, with a 1.5 safety factor applied at motor selection
\cite{siegwart2011mobilerobots};
TEC1-12706 modules of 40\,mm $\times$ 40\,mm \cite{tec12706ds} in a
2$\times$2 bench array with 2\,m/s forced air on the cold side
\cite{zhao2022teccondenser}; a 3.6\,L tank sized for one cleaning event at
1.8\,L/m$^2$ over 2\,m$^2$; and an ESP32 electronics bay with INA219 current
sensing per the prototype test plan.

\section{Drawing register}
\begin{table}[h]
\centering
\caption{Drawing register. Sheets D-101 to D-107 dated 2026-06-12; the
rail-drive sheets D-108 to D-110 dated 2026-06-13. Units mm, general
tolerance ISO 2768-m.}
\label{tab:register}
\small
\begin{tabular}{lllll}
\toprule
No. & Title & Material & Scale & Files (\fpath{cad/}) \\
\midrule
D-101 & Chassis base plate & Al 6061-T6, 3\,mm & 1:2 A3 &
  \fpath{D-101\_chassis\_base\_plate.scr/.dxf} \\
D-102 & Motor mount bracket ($\times$2) & POM & 1:1 A4 &
  \fpath{D-102\_motor\_mount\_bracket.scr/.dxf} \\
D-103 & Brush assembly & mixed (see sheet) & 1:2 A3 &
  \fpath{D-103\_brush\_assembly.scr/.dxf} \\
D-104 & TEC cold-plate sandwich & Al 6061 + Al extrusion & 1:2 A3 &
  \fpath{D-104\_tec\_coldplate\_sandwich.scr/.dxf} \\
D-105 & Water tank 3.6\,L + pump & HDPE, 2\,mm & 1:2 A3 &
  \fpath{D-105\_water\_tank\_pump.scr/.dxf} \\
D-106 & Electronics bay & acrylic, 3\,mm & 1:1 A4 &
  \fpath{D-106\_electronics\_bay.scr/.dxf} \\
D-107 & General assembly & --- & 1:2 A3 &
  \fpath{D-107\_general\_assembly.scr/.dxf} \\
\addlinespace
D-108 & Rail assembly (2$\times$2.4\,m) & Al 6063 2040 V-slot & 1:5 A3 &
  see \fpath{cad/3d/rail\_and\_panel\_3d.scr} \\
D-109 & Rack-and-pinion drive & steel rack + pinion, $\varnothing$30 & 1:2 A3 &
  see \fpath{cad/3d/pvbot\_3d\_assembly.scr} \\
D-110 & Full-width brush ($0.96$\,m) & microfibre over $\varnothing$50 core & 1:5 A3 &
  see \fpath{cad/3d/pvbot\_3d\_assembly.scr} \\
\bottomrule
\end{tabular}
\end{table}

The three rail-drive sheets D-108 to D-110 are carried in the 3D model in
\fpath{cad/3d/} rather than as separate 2D scripts: \fpath{rail\_and\_panel\_3d.scr}
builds the rails, rack, posts, and end stops (D-108), and
\fpath{pvbot\_3d\_assembly.scr} builds the rack-and-pinion drive (D-109) and the
full-width brush (D-110) as part of the carriage. \fpath{cad/3d/README\_3D.md}
lists the exact AutoCAD steps and the workshop dimensions chosen for the model.

Note on numbering: documents from the earlier Benban skid phase used D-101 to
D-107 for the fixed AWG skid. After the April 2026 pivot the robot is the
deliverable, so this register redefines the same numbers for the robot parts;
the skid register in the older documents is superseded.

\subsection{Design revision: free crawler to rail-guided carriage}
\label{sec:railrev}
The earlier drawings describe a free differential-drive crawler that drove a
0.25\,m brush over the panel in a boustrophedon (lawn-mower) path. The Gazebo
test of that crawler on the $24^\circ$ glass showed it could not hold a
heading: with the panel friction the report assumes ($\mu=0.6$ glass under
$\mu=0.9$ rubber), the wheels broke traction during every in-place yaw, so the
robot could not complete the turns the boustrophedon path needs and drifted
down the slope. The robot was therefore changed to a rail-guided carriage. Two
rails run down the fall-line, the carriage rides them on V-groove wheels, and a
rack-and-pinion drive moves it up and down; there is no yaw to lose. The brush
now spans the full panel width, so one sweep down and back cleans the whole
panel with no turning. The changed quantities are in
\fpath{output/rail\_guided\_design.json}:

\begin{itemize}\setlength\itemsep{1pt}
  \item Brush: 0.25\,m swath $\rightarrow$ full-width 0.96\,m at $\varnothing$50\,mm.
  \item Path: 5-pass boustrophedon (10.9\,m, 230\,s) $\rightarrow$ single
        fall-line sweep down and up, 66.7\,s per panel.
  \item Drive: two free traction wheels $\rightarrow$ rack-and-pinion on two
        rails, 0.84\,N\,m at the $\varnothing$30\,mm pinion at 38\,rpm.
  \item Mass: 6.0\,kg crawler $\rightarrow$ 8.0\,kg carriage (rails wheels,
        side-frames, drive, longer brush).
\end{itemize}

The TEC water module, the 3.6\,L tank, the spray system, the power chain, and
the sensors are unchanged. The buy-list for the rail hardware, drive motor with
encoder, rack, pinion, and the full-width brush is
\fpath{MARKET\_PARTS\_BUYLIST.md}, with live verified vendor links. The 3D
model of the revised robot on its rails is \fpath{cad/3d/} (see
Section~\ref{sec:register3d}).

\section{How to open the drawings}

\subsection{AutoCAD (.scr scripts)}
Each \fpath{.scr} file is a plain-text command script that draws the part at
1:1 in millimetres, with layers, dimensions, title block, and notes.
\begin{enumerate}
  \item Open AutoCAD and start a new drawing from the \fpath{acadiso.dwt}
        template (metric).
  \item Type \texttt{SCRIPT}, press Enter, and browse to the \fpath{.scr}
        file (or drag the file into the drawing window).
  \item The script sets its own layers and dimension variables, draws the
        views, and ends with \texttt{ZOOM E}. Save as \fpath{.dwg} if needed.
\end{enumerate}
Blank lines inside the scripts are deliberate Enter keystrokes; do not strip
them if editing. The scripts only use core commands (\texttt{LINE},
\texttt{CIRCLE}, \texttt{RECTANG}, \texttt{DIMLINEAR}, \texttt{-TEXT},
\texttt{-LAYER}), so they also run in LT and in older releases.

\subsection{DXF files}
The \fpath{.dxf} files (AutoCAD R2010 format) open directly with
\texttt{OPEN} in AutoCAD, and in any DXF viewer (LibreCAD, QCAD, FreeCAD,
DWG TrueView). They carry the same geometry, layers, and dimensions as the
scripts and were verified with the \fpath{ezdxf} auditor (0 errors on all
seven files). To regenerate after a dimension change:
\texttt{python scripts/make\_cad\_dxf.py}.

\subsection{Gazebo model}
\fpath{gazebo/pvbot\_description/} holds \fpath{pvbot.urdf.xacro} (robot,
primitive geometry only), \fpath{panel\_world.sdf} (2\,m$^2$ panel at
$24^\circ$ on a ground plane), and \fpath{README\_GAZEBO.md} with exact
spawn commands for Gazebo Harmonic (\texttt{gz sim} + \texttt{ros\_gz}) and
Gazebo Classic 11 (\texttt{gazebo\_ros}). In short: convert with
\texttt{xacro}, launch the world, then spawn with
\texttt{ros2 run ros\_gz\_sim create} (Harmonic) or
\texttt{ros2 run gazebo\_ros spawn\_entity.py} (Classic), and drive at the
0.06\,m/s cleaning speed over \texttt{/cmd\_vel}.

\subsection{3D model (\fpath{cad/3d/})}
\label{sec:register3d}
The 3D model of the revised rail-guided robot is a set of AutoCAD scripts that
build native 3D solids. \fpath{pvbot\_3d\_assembly.scr} builds the carriage:
chassis plate, TEC AWG module (cold plate, four tiles, finned hot sink, two
fans), 3.6\,L tank, full-width 0.96\,m brush, spray bar with nine nozzles,
side-frames with V-groove wheels, and the rack-and-pinion drive.
\fpath{rail\_and\_panel\_3d.scr} builds the $24^\circ$ panel on posts, the two
2.4\,m V-slot rails, the rack strips, and the end stops.
\fpath{scene\_3d.scr} nests the two so the carriage sits on the rails in one
drawing. Each subsystem is on its own colour-coded layer. To open: start a
metric drawing, type \texttt{SCRIPT}, pick the file, then \texttt{ZOOM E},
\texttt{-VIEW SWISO}, and \texttt{VSCURRENT} Realistic. A faceted-mesh copy,
\fpath{cad/3d/pvbot\_3d.dxf} (regenerated by
\fpath{cad/3d/make\_pvbot\_3d\_dxf.py}), opens in any DXF viewer.
\fpath{cad/3d/README\_3D.md} carries the full steps and the workshop
dimensions chosen for the model. Sizes match \fpath{pvbot.urdf} and the frozen
rail-guided values.

\section{Design rationale by part}

\paragraph{D-101 chassis base plate.}
The plate is the 400\,mm $\times$ 300\,mm footprint of the URDF chassis box,
cut from 3\,mm Al 6061-T6. Thickness is set by stiffness, not strength: the
worst static load is the full 6\,kg robot reacting through the four wheel and
caster points ($\approx$59\,N total), and bending stress in a 3\,mm plate at
this load is below 5\,MPa against a 240\,MPa yield, while the 3\,mm section
keeps deflection under the brackets small enough that wheel camber does not
change measurably under load. The traction reaction (32\,N along the deck on
a $24^\circ$ panel) passes through the eight M4 bracket bolts in shear, each
loaded to a few newtons. Wheels run in two 80\,mm $\times$ 26\,mm slots so
the axle sits low (14\,mm above the deck) and the chassis keeps the 15\,mm
ground clearance of the URDF model. All bays and hole groups are tabulated
on the sheet (hole table A--F).

\paragraph{D-102 motor mount bracket ($\times$2).}
The bracket carries the geared DC motor whose rating follows from the frozen
traction calculation: 1.37\,N\,m total, 0.686\,N\,m per wheel, 1.03\,N\,m
selected with the 1.5 safety factor \cite{siegwart2011mobilerobots}. At
stall the bore reacts 1.03\,N\,m as a force couple across the 28\,mm
$\times$ 44\,mm M4 base pattern, about 23\,N per bolt, two orders below M4
single-shear capacity. POM is chosen because it machines cleanly, damps gear
noise, and will not fret against the anodised plate. The $\varnothing$25 H8
clamp bore fits the JGA25-370 gearmotor body ($\varnothing$24.4); a 2\,mm
slit with two M3 screws clamps it without a keyway.

\paragraph{D-103/D-110 brush assembly.}
After the rail revision the roller spans the full 0.96\,m panel width (sheet
D-110), so a single fall-line sweep cleans the panel with no turning; the
earlier 0.25\,m swath and its 5-pass boustrophedon path are superseded
(Section~\ref{sec:railrev}). The $\varnothing$50\,mm roller core matches the
URDF \texttt{brush\_link} cylinder. A microfibre sleeve over a $\varnothing$25
PVC core keeps the roller soft against the glass; the longer roller raises the
brush-assembly mass to about 1.0\,kg. The $\varnothing$8 steel shaft runs in
two 608ZZ bearings in POM end housings, sized so bearing load stays
negligible. The spray bar sits behind the brush with nine $\varnothing$1.0
nozzles at 100\,mm pitch across the width, angled ahead of the contact line so
water lands just before the brush; the low-pressure diaphragm pump (D-105)
delivers the 1.8\,L/m$^2$ budget without atomising.

\paragraph{D-104 TEC cold-plate sandwich.}
Four TEC1-12706 tiles \cite{tec12706ds} sit in 40\,mm $\times$ 40\,mm
$\times$ 0.5\,mm pockets at 85\,mm pitch in a 170\,mm $\times$ 170\,mm
$\times$ 3\,mm Al 6061 cold plate, the 2$\times$2 bench layout of the
prototype chapter. The pocket locates each tile and thins the conduction
path; thermal paste and M4 clamps with nylon shoulder washers (1.2\,N\,m)
give full-face contact. The hot side is a 160\,mm $\times$ 160\,mm field of
four 80\,mm $\times$ 80\,mm $\times$ 25\,mm extruded sinks with two 80\,mm
fans blowing through the fins, the arrangement assumed by the 0.15\,K/W
hot-side resistance in the operating-point solver
\cite{zhao2022teccondenser,lin2022heatsink}. A separate ducted fan gives the
2\,m/s cold-side face velocity used in the condensation correlation. The
plate mounts tilted $15^\circ$ so condensate runs to a 12\,mm $\times$
12\,mm gutter and a $\varnothing$8 drain into the tank. The controller
floors the plate at $+1\,^\circ$C as frost protection.

\paragraph{D-105 water tank and pump.}
Tank capacity is exactly one cleaning event: 1.8\,L/m$^2$ over 2\,m$^2$
gives 3.6\,L, achieved with internal dimensions 166\,mm $\times$ 226\,mm
$\times$ 96\,mm in 2\,mm welded HDPE. The flat, low shape keeps the full
tank's centre of gravity about 70\,mm above the deck, which matters on a
$24^\circ$ slope (Section~\ref{sec:ga}). Lugs match the four $\varnothing$5
strap holes on D-101. Fittings: $\varnothing$40 filler, $\varnothing$8
condensate inlet from the D-104 drain, $\varnothing$8 outlet barb to the
12\,V diaphragm pump (90\,mm $\times$ 40\,mm $\times$ 35\,mm, feet on the
two M4 plate holes), reed level switch in a $\varnothing$16 boss, and a
$\varnothing$4 vent.

\paragraph{D-106 electronics bay.}
A 150\,mm $\times$ 100\,mm $\times$ 3\,mm acrylic deck on M3 standoffs
carries the ESP32 DevKitC, the INA219 current monitor called for by the test
plan, the dual motor driver, and the buck converter, with the small 12\,V
battery strapped along the front edge. Acrylic is non-conductive, laser cuts
in one pass, and lets the wiring be inspected through the matching vented
cover. Standoff patterns follow the common module footprints (48\,mm
$\times$ 22\,mm for the DevKit, 37\,mm $\times$ 37\,mm for the driver) and
are dimensioned on the sheet, so substituting boards only means redrilling
the deck, not the chassis.

\paragraph{D-107 general assembly.}
The GA fixes the envelope at 460\,mm $\times$ 300\,mm $\times$ 155\,mm
(195\,mm when the $15^\circ$ harvest shim is fitted under the TEC sandwich)
and prints the assembly order: plate, brackets and motors, wheels, casters,
tank and pump, TEC sandwich on the tank lid, electronics bay, brush
assembly, wiring last. A tipping check justifies the layout: the combined
centre of gravity sits about 80\,mm above the panel surface and 140\,mm
inside the rear wheel contact, so on a $24^\circ$ slope the overturning
lever ($80 \times \tan 24^\circ \approx 36$\,mm) is a quarter of the
restoring lever and the robot cannot tip before it slides; sliding itself is
held by the wheel friction included in the 32\,N traction figure.

\section{Designer's-choice assumptions}
The frozen calculation fixes masses, torques, the swath, tile size, and tank
volume, but not every workshop dimension. Table~\ref{tab:assumptions} lists
the values chosen here so they can be revisited without touching the frozen
analysis.

\begin{table}[h]
\centering
\caption{Dimensions chosen in this package (not fixed by the frozen
calculations).}
\label{tab:assumptions}
\small
\begin{tabular}{p{0.42\textwidth}p{0.23\textwidth}p{0.27\textwidth}}
\toprule
Assumption & Value & Basis \\
\midrule
Chassis plate size and thickness & 400$\times$300$\times$3\,mm & URDF box
footprint; stiffness \\
Wheel slots in plate & 2$\times$ 80$\times$26\,mm & wheel $\varnothing$60
$\times$20 + clearance \\
Axle height above deck & 14\,mm & motor $\varnothing$24.4 clearance; URDF
uses 12\,mm (rounded box model) \\
Drive motor envelope & JGA25-370, $\varnothing$24.4$\times$62\,mm & common
$\geq$1\,N\,m geared DC class \\
Bracket base / bore & 44$\times$60$\times$6\,mm / $\varnothing$25\,H8 &
clamp fit on motor body \\
Roller core / shaft / bearings & $\varnothing$25 PVC / $\varnothing$8\,h6 /
608ZZ & lightest stock parts \\
Spray nozzles & 5$\times$ $\varnothing$1.0 at 50\,mm pitch & even wetting of
250\,mm swath \\
Cold plate / TEC pitch / pocket & 170$\times$170$\times$3\,mm / 85\,mm /
0.5\,mm deep & 2$\times$2 layout, sink field width \\
Hot sinks and fans & 4$\times$ 80$\times$80$\times$25\,mm + 2$\times$80\,mm
fans & matches 0.15\,K/W solver value \\
Gutter and drain & 12$\times$12\,mm, $\varnothing$8 & gravity drain at
$15^\circ$ tilt \\
Tank internal size, wall & 166$\times$226$\times$96\,mm, 2\,mm HDPE &
3.60\,L working volume \\
Pump & 12\,V diaphragm, 90$\times$40$\times$35\,mm & prototype BOM class \\
Electronics deck and cover & acrylic 150$\times$100$\times$3\,mm & module
footprints + battery strip \\
Overall envelope & 460$\times$300$\times$155\,mm & GA stack-up \\
\bottomrule
\end{tabular}
\end{table}

\section{Mass budget}
\label{sec:mass}
The rail-guided carriage budget closes at the frozen 8.00\,kg with the tank at
mean fill (1.8\,kg of water). Against the earlier 6.0\,kg crawler, the
revision drops the two free traction wheels and the casters and adds the
carriage side-frames, the eight V-groove guide wheels, the two-rail
rack-and-pinion drive, and the heavier full-width brush (Table~\ref{tab:mass}).
The Gazebo URDF, \fpath{pvbot.urdf}, is the 6.0\,kg differential-drive
simulation skeleton and is not re-massed here; it keeps the validated inertias
used in the traction test that drove the revision.

\begin{table}[h]
\centering
\caption{Rail-guided carriage mass budget at mean tank fill. Closure 8.00\,kg.}
\label{tab:mass}
\small
\begin{tabular}{llr}
\toprule
Item & Drawing & Mass (kg) \\
\midrule
Carriage base plate & D-101 & 1.10 \\
TEC sandwich (tiles, plate, sinks, fans) & D-104 & 1.05 \\
Tank + pump, dry & D-105 & 0.40 \\
Water, mean fill (1.8\,L) & --- & 1.80 \\
Electronics bay incl.\ battery & D-106 & 0.37 \\
Full-width brush ($0.96$\,m, $\varnothing$50) & D-103/D-110 & 1.00 \\
Carriage side-frames (2) & D-108 & 0.80 \\
V-groove guide wheels ($\times$8) & D-108 & 0.40 \\
Rack-and-pinion drive (2 gearmotors + pinions) & D-109 & 0.46 \\
IMU, sensors, cabling & --- & 0.62 \\
\midrule
\textbf{Total} & & \textbf{8.00} \\
\bottomrule
\end{tabular}
\end{table}

Two consequences follow. First, water is about 23\,\% of the moving mass, so
the drive margin grows as the tank empties during cleaning; the full-mass case
is the worst case for the rail drive. Second, the drive sizing in
\fpath{rail\_guided\_design.json} is 37.2\,N at the rack and 0.84\,N\,m at the
$\varnothing$30\,mm pinion at full carriage mass on the $24^\circ$ panel;
splitting the load across two driven rails halves the per-motor torque, so a
GA25-370-class gearmotor with encoder carries it inside a 1.5 safety factor.

\section{Gazebo model and URDF correspondence}
\label{sec:ga}
\fpath{pvbot.urdf.xacro} keeps every value of the architecture document's
skeleton: \texttt{base\_link} box 0.40$\times$0.30$\times$0.12\,m at
5.0\,kg with inertia diag(0.0435, 0.0727, 0.1042)\,kg\,m$^2$; wheels of
0.03\,m radius and 0.15\,kg at $(-0.12, \pm 0.13, -0.045)$ on a 0.26\,m
track; the caster as a frictionless sphere; the brush as a continuous joint
at $(0.22, 0, -0.05)$. This package adds what the skeleton omitted: visual
primitives with materials, caster and brush inertials (sphere and cylinder
formulas: $8.1\times10^{-6}$ and diag($3.4\times10^{-3}$,
$3.4\times10^{-3}$, $4.8\times10^{-4}$)\,kg\,m$^2$), a wheel effort limit of
1.03\,N\,m, and a differential-drive plugin for each simulator generation.
The world file \fpath{panel\_world.sdf} is the validation scene: a 2.0\,m
$\times$ 1.0\,m panel (2\,m$^2$, the frozen cleaned area) pitched at
$24^\circ$, with glass-like friction ($\mu = 0.6$) under wheel rubber at
$\mu = 0.9$. Driving the model at 0.05\,m/s with full tank mass reproduces
the torque sizing case of the report in simulation.

\section{Traceability of frozen values}
Every number frozen by the verified calculation engine appears somewhere a
manufacturer or simulator will read it. Table~\ref{tab:trace} maps each
frozen value to its location in the drawings and the model, so a change in
the analysis can be propagated by inspection.

\begin{table}[h]
\centering
\caption{Frozen design values and where each one lands in the package.}
\label{tab:trace}
\small
\begin{tabular}{p{0.30\textwidth}p{0.16\textwidth}p{0.44\textwidth}}
\toprule
Frozen value & Source & Appears in \\
\midrule
Robot mass 6.0\,kg & results JSON & D-107 item list; Table~\ref{tab:mass};
URDF mass sum \\
Wheel radius 0.03\,m & results JSON & D-107 wheel $\varnothing$60; URDF
wheel cylinders \\
Track 0.26\,m & URDF skeleton & D-101 axle note; D-107 plan dim; URDF
joint origins and plugin \\
Resistive force 32.0\,N & \fpath{mechanics.F\_resist} & D-107 design-basis
note; rationale D-101/D-102 \\
Traction torque 1.37\,N\,m & \fpath{mechanics.T\_traction} & D-107 note;
D-102 rationale \\
Motor rating 1.03\,N\,m (SF 1.5) & report Ch.~6 & D-102 rationale; URDF
joint effort limit \\
Brush swath 0.25\,m & report Ch.~6 & D-103 roller face; D-107 plan dim;
URDF brush length \\
Panel tilt $24^\circ$, area 2\,m$^2$ & report Ch.~1, 6 & D-107 design
basis; \fpath{panel\_world.sdf} pose and panel size \\
Cleaning speed 0.05\,m/s & report Ch.~6 & README spawn commands
(\texttt{cmd\_vel}) \\
TEC1-12706, 40$\times$40\,mm, 4 tiles & report Ch.~2, 8 & D-104 pockets and
section \\
Cold-side air 2\,m/s & report Eq.~(3.7) & D-104 note 2 \\
Hot-side $R_{th}=0.15$\,K/W & report Ch.~4 & D-104 sink/fan arrangement \\
Frost floor $+1\,^\circ$C & report Eq.~(3.13) & D-104 note 4 \\
Tank 3.6\,L (1.8\,L/m$^2$) & \fpath{cleaning.water\_per\_event} & D-105
internal volume and level line \\
ESP32 + INA219 instrumentation & report Ch.~7, 9 & D-106 standoff layout \\
\bottomrule
\end{tabular}
\end{table}

\section{Manufacturing notes and fasteners}
All plate parts (D-101 chassis, D-106 deck and cover, D-104 fan plate) are
single-pass laser or waterjet cuts; only the D-102 brackets and the D-104
cold-plate pockets need milling. The tank is welded from HDPE sheet, which a
local plastics shop can do from the D-105 blank dimensions; a rotomoulded or
off-the-shelf tank of the same internal volume is an acceptable substitute
if the lug positions are reproduced. Anodising the chassis plate is optional
but recommended on the deck face that carries the POM brackets.

Hole-to-fastener mapping across the set: $\varnothing$3.2 takes M3
(electronics standoffs, motor face and clamp screws), $\varnothing$4.5
takes M4 (brackets, casters, pump, brush arms, TEC clamps), and
$\varnothing$5 takes the M5 tank strap bolts. Countersinks appear only at
the caster holes (M4 flat head, so the heads clear the panel surface at full
suspension compression). The complete fastener kit is: M3$\times$8 (12),
M3$\times$6 (8), M3 brass standoffs 10\,mm (4) and 35\,mm (4), M4$\times$12
(18), M4$\times$16 flat head (4), M5$\times$16 (4), plus nyloc nuts
throughout; total hardware mass is inside the 0.10\,kg miscellaneous
allowance distributed across the item masses of Table~\ref{tab:mass}.

Tolerances follow ISO 2768-m everywhere except three fits called out on the
sheets: the D-102 motor bore ($\varnothing$25\,H8, reamed), the D-103
bearing seats (608ZZ outer ring, $\varnothing$22\,H7), and the D-101 motor
bracket hole group (positioned $\pm$0.2\,mm to the axle centreline so the
two wheels stay parallel within the differential-drive controller's
assumption).

\section{File manifest}
\begin{itemize}\setlength\itemsep{0pt}
  \item \fpath{cad/D-101\dots D-107 (.scr + .dxf)}: drawing set (14 files)
  \item \fpath{scripts/make\_cad\_dxf.py}: DXF generator (single
        dimension source for the DXF set)
  \item \fpath{gazebo/pvbot\_description/pvbot.urdf.xacro}: robot model
  \item \fpath{gazebo/pvbot\_description/panel\_world.sdf}: $24^\circ$
        panel world
  \item \fpath{gazebo/pvbot\_description/README\_GAZEBO.md}: spawn
        commands (Harmonic and Classic)
  \item \fpath{CAD\_DESIGN\_PACKAGE.tex / .pdf}: this document
\end{itemize}

\bibliographystyle{ieeetr}
\bibliography{references}

\end{document}
